\section{线性发散的简单检验}\label{sec:test}

为研究准 PDF 算符的重正化性质，需要在多个不同格点间距下计算矩阵元。为确保数据可用于精细的高精度分析，我们先检验这些数据是否近似呈现出微扰论预言的线性发散，特别要证明线性发散项的 $z$ 依赖是线性的。

为此，我们先从裸矩阵元中提取因子 $e^{-\delta \bar{m} z}$，以检验式~(\ref{eq:delta_m}) 中 $\delta \bar{m}$ 的 $1/a$ 依赖。基于式~(\ref{eq:Reoperator})、(\ref{eq:delta_m}) 和~(\ref{eq:alpha_s})，对裸矩阵元 $\mathcal{M}$ 取自然对数可得
\bea\label{eq:logM_s}
\ln \mathcal{M}(z,a) = \frac{e(z)}{a \ln[a \Lambda_{\rm QCD}]} + g(z),
\eea
其中第一项线性发散，$g(z)$ 是剩余项。这里我们忽略了对数发散和离散化误差，因为它们小得多，对该线性发散初步检验几乎没有影响。

\begin{figure}[tbp]
\centering
\includegraphics[width=8cm]{images/fitting_RIMOM_overlap.pdf}
\includegraphics[width=8cm]{images/ez_RIMOM_overlap.pdf}
\caption{
上图：用式~(\ref{eq:logM_s}) 拟合无 HYP 涂抹、overlap 作用下离壳夸克态中裸 $O_{\gamma_t}(z)$ 矩阵元。蓝点为插值数据，彩色曲线为每个 $z$ 的拟合曲线。$\Lambda_{\rm QCD}$ 固定为 0.2 GeV。
下图：$e(z)$ 随 $z$ 的变化。蓝点是上图中拟合得到的参数 $e(z)$，蓝色曲线是 $e(z)$ 的无偏线性拟合。}
\label{fig:testlindiv_RIMOM_overlap}
\end{figure}

\begin{figure}[tbp]
\centering
\includegraphics[width=8cm]{images/fitting_RIMOM_clover.pdf}
\includegraphics[width=8cm]{images/ez_RIMOM_clover.pdf}
\caption{同图~\ref{fig:testlindiv_RIMOM_overlap}，但价夸克为 clover 作用量。}
\label{fig:testlindiv_RIMOM_clover}
\end{figure}

我们用式~(\ref{eq:logM_s}) 拟合裸矩阵元随 $a$ 的变化。将 $e(z)$ 与 $g(z)$ 视为 $z$ 的未知函数，在每个 $z$ 值处分别进行拟合，把它们当作自由参数。需要对数据 $\ln \mathcal{M}(z,a)$ 关于 $z$ 作线性插值。具体步骤如下：

1. 用邻近的两个数据点作线性插值，得到 $z=0.06\times n$ fm $(n=3,4,5,...,16)$ 处的中心值与不确定度。

2. 对每个 $z$，拟合其对 $a$ 的依赖，将 $e(z)$、$g(z)$ 作为拟合参数。在这个初步测试中，我们把 $\Lambda_{\rm QCD}$ 固定为 0.2 GeV，这是一个合理的起始猜测~\cite{Lepage:1992xa}。

3. 绘出拟合得到的 $e(z)$ 随 $z$ 的变化，以考察 $e(z)$ 是否具有线性的 $z$ 依赖。

如果第 2 步得到合理的拟合，那么就能看到发散近似呈 $1/a$ 依赖；第 3 步则可检验线性的 $z$ 依赖。无 HYP 涂抹时，价 overlap 与 clover 作用下离壳夸克态裸 $O_{\gamma_t}(z)$ 矩阵元的拟合结果分别示于图~\ref{fig:testlindiv_RIMOM_overlap} 与图~\ref{fig:testlindiv_RIMOM_clover}。共有八档不同的格点间距，$z$ 的范围取 0.18~fm 至 0.96~fm，其间分析了 14 个不同的 $z$ 值。

两种情形下的 $1/a$ 依赖都拟合得很好；不过由于数据精度很高，卡方值较大，但这不是本初步阶段关心的问题。发散系数的 $z$ 依赖很好地展现出线性特征，且大致经过 $z=0$ 处的零点。这有力地表明线性发散大体遵循微扰论的预言，使我们有信心在后续几节中进行更精细的分析。
